\documentclass[12pt,a4paper]{article}
\usepackage[utf-8]{inputenc}
\usepackage[margin=1in]{geometry}
\usepackage{graphicx}
\usepackage{amsmath}
\usepackage{amssymb}
\usepackage{xcolor}
\usepackage{listings}
\usepackage{hyperref}
\usepackage{fancyhdr}
\usepackage{tcolorbox}
\usepackage{enumerate}
\usepackage{booktabs}
\usepackage{array}
\usepackage{textcomp}

% Fix for listings package
\lstset{
    language=Python,
    basicstyle=\ttfamily\footnotesize,
    keywordstyle=\color{blue}\bfseries,
    commentstyle=\color{gray},
    stringstyle=\color{red},
    breaklines=true,
    showstringspaces=false,
    tabsize=4,
    frame=lines,
    rulecolor=\color{black},
    numbers=left,
    numberstyle=\tiny,
    stepnumber=1
}

% Header and footer
\pagestyle{fancy}
\fancyhf{}
\lhead{MSVTNet Project}
\rhead{Interview Preparation}
\cfoot{\thepage}

\title{\textbf{\Large MSVTNet Project Explanation}\\
\textbf{Brain Signal Classification using Vision Transformers}\\
\normalsize Interview Guide \& Data Validation Role}
\author{Data Quality \& Validation Specialist}
\date{\today}

\begin{document}

\maketitle

\tableofcontents
\newpage

% ============================================================================
\section{Project Overview}
% ============================================================================

\subsection{What is MSVTNet?}

MSVTNet stands for \textbf{Multi-Scale Vision Transformer Neural Network}. It is an artificial intelligence system designed to:

\begin{itemize}
    \item Read \textbf{brain signals} (EEG - Electroencephalography)
    \item Recognize what \textbf{movement} a person is \textbf{thinking about}
    \item Classify into 4 movement types: Left Hand, Right Hand, Feet, or Tongue
    \item Achieve \textbf{94\% plus accuracy} on real brain signals
\end{itemize}

\subsection{Real-World Application}

Imagine someone who is paralyzed. Using this technology:

\begin{tcolorbox}[colback=blue!10,colframe=blue!50!black,title=Use Case]
Person thinks ``move left hand'' $\rightarrow$ System recognizes it $\rightarrow$ Robotic arm moves \textbf{automatically}
\end{tcolorbox}

\subsection{Project Technology Stack}

\begin{table}[h]
\centering
\begin{tabular}{|l|l|l|}
\hline
\textbf{Component} & \textbf{Technology} & \textbf{Purpose} \\
\hline
Language & Python 3.10 & Programming \\
Deep Learning & PyTorch 2.0.1 & AI Model \\
GPU Support & CUDA 11.8 & Fast Processing \\
Data Processing & NumPy, Pandas & Data Validation \\
Hardware & NVIDIA RTX 3090 & GPU Acceleration \\
\hline
\end{tabular}
\end{table}

% ============================================================================
\section{Your Role: Data Quality and Validation Specialist}
% ============================================================================

\subsection{Responsibilities}

Your primary responsibility is to:

\begin{enumerate}
    \item \textbf{Check for missing values} (NULL values) in datasets
    \item \textbf{Validate data integrity} across all files
    \item \textbf{Identify anomalies} in brain signal data
    \item \textbf{Generate reports} on data quality
    \item \textbf{Ensure data consistency} before model training
\end{enumerate}

\subsection{Tools You Use}

\begin{tcolorbox}[colback=yellow!10,colframe=orange!50!black,title=Your Toolkit]
\begin{itemize}
    \item \textbf{NumPy}: For numerical array operations
    \item \textbf{Pandas}: For data frame manipulation and analysis
    \item \textbf{Python}: For scripting and automation
\end{itemize}
\end{tcolorbox}

% ============================================================================
\section{Project Pipeline: Step-by-Step Explanation}
% ============================================================================

\subsection{Step 1: Data Collection}

\subsubsection{Non-Technical Explanation}

Think of this as \textbf{recording brain activity}:

\begin{tcolorbox}[colback=green!10,colframe=green!50!black,title=Recording Process]
\begin{enumerate}
    \item Put an \textbf{EEG headset} on a person (like a cap with 22 electrodes)
    \item Ask them to \textbf{THINK about moving}:
    \begin{itemize}
        \item Think about moving your LEFT HAND: Record signal
        \item Think about moving your RIGHT HAND: Record signal
        \item Think about moving your FEET: Record signal
        \item Think about moving your TONGUE: Record signal
    \end{itemize}
    \item Repeat 40-50 times per movement type
    \item Save as GDF files (raw brain signal format)
\end{enumerate}
\end{tcolorbox}

\subsubsection{Technical Explanation}

\begin{lstlisting}[caption={Raw Data Format}]
Raw Brain Signals (GDF file):
- 22 channels (electrodes)
- 250 Hz sampling rate (250 samples per second)
- 3.5 seconds per trial
- Total samples per trial: 250 x 3.5 = 875 data points

Data shape: (number_of_trials, 22_channels, 875_timepoints)
Example: (40 trials, 22 channels, 875 timepoints)
\end{lstlisting}

\subsection{Step 2: Data Preprocessing (Your Quality Check Point!)}

\subsubsection{Non-Technical Explanation}

Raw brain signals are \textbf{MESSY}. They contain:

\begin{itemize}
    \item \textbf{Noise}: Electrical interference from the environment
    \item \textbf{Artifacts}: Muscle movements, eye blinks
    \item \textbf{Inconsistencies}: Different signal scales between people
\end{itemize}

\textbf{Preprocessing = Cleaning the data}

\subsubsection{Technical Explanation}

\begin{tcolorbox}[colback=cyan!10,colframe=cyan!50!black,title=Preprocessing Steps]
\begin{enumerate}
    \item \textbf{Bandpass Filter (4-40 Hz)}: Keep only brain signal frequencies
    \item \textbf{Common Average Reference}: Remove common noise
    \item \textbf{Standardization}: Normalize to mean=0, std=1
    \item \textbf{Save to NPZ}: Compressed format for storage
\end{enumerate}
\end{tcolorbox}

\subsubsection{Your Quality Check Role}

\begin{lstlisting}[caption={Check Data After Preprocessing}]
import numpy as np
import pandas as pd
import glob

# Load preprocessed data
data_path = "path/to/preprocessed/data"
npz_files = glob.glob(f"{data_path}/*.npz")

for file in npz_files:
    # Load NPZ file
    data = np.load(file)
    
    # Extract data and labels
    signals = data['data']
    labels = data['labels']
    
    # Check 1: Any NULL/NaN values?
    print("Checking for NULL values...")
    nan_count_signals = np.isnan(signals).sum()
    nan_count_labels = np.isnan(labels).sum()
    
    print(f"NaN in signals: {nan_count_signals}")
    print(f"NaN in labels: {nan_count_labels}")
    
    # Check 2: Signal range
    print(f"Signal range: {signals.min()}, {signals.max()}")
    
    # Check 3: Label distribution
    unique_labels, counts = np.unique(labels, return_counts=True)
    print(f"Class distribution: {dict(zip(unique_labels, counts))}")
\end{lstlisting}

\subsection{Step 3: Data Loading for Training}

\subsubsection{Non-Technical Explanation}

Before the AI trains, data must be:

\begin{enumerate}
    \item \textbf{Split}: 70\% for training, 30\% for testing
    \item \textbf{Batched}: Grouped into chunks of 32 samples
    \item \textbf{Prepared}: Converted to PyTorch format
\end{enumerate}

\subsubsection{Technical Explanation}

\begin{lstlisting}[caption={Data Validation Before Training}]
import torch
from torch.utils.data import DataLoader, Dataset

# Create custom dataset
class BCIDataset(Dataset):
    def __init__(self, data, labels):
        self.data = data
        self.labels = labels
        
        # Validate during initialization
        assert self.data.shape[0] == self.labels.shape[0], \
            "Data and labels must have same length"
        
        assert not np.isnan(self.data).any(), \
            "Found NaN values in data!"

# Create DataLoader
train_loader = DataLoader(train_dataset, 
                         batch_size=32, shuffle=True)

# Your role: Verify data quality
for batch_data, batch_labels in train_loader:
    if torch.isnan(batch_data).any():
        print("ERROR: NaN in batch data!")
    else:
        print("Batch OK")
    break
\end{lstlisting}

\subsection{Step 4: Model Architecture}

\subsubsection{Non-Technical Explanation}

The MSVTNet model has 3 main blocks:

\begin{tcolorbox}[colback=magenta!10,colframe=magenta!50!black,title=Model Structure]

\textbf{BLOCK 1: Multi-Scale Extraction (MSST)}
\begin{itemize}
    \item Looks at brain signals with 3 different zoom levels
    \item Window sizes: 3, 5, 7
    \item Output: 3 different feature maps
\end{itemize}

\textbf{BLOCK 2: Feature Combination (CSGT - Transformer)}
\begin{itemize}
    \item Takes 3 feature maps from Block 1
    \item Figures out which features are related
    \item Uses attention mechanism (selective focus)
    \item Output: 1 combined feature representation
\end{itemize}

\textbf{BLOCK 3: Classification (Classifier Head)}
\begin{itemize}
    \item Takes combined features
    \item Predicts: Left Hand / Right Hand / Feet / Tongue
    \item Output: 1 of 4 predictions
\end{itemize}

\end{tcolorbox}

\subsubsection{Technical Explanation}

\begin{lstlisting}[caption={Model Architecture}]
import torch.nn as nn

class MSVTNet(nn.Module):
    def __init__(self, num_channels=22, num_classes=4):
        super(MSVTNet, self).__init__()
        
        # BLOCK 1: Multi-Scale Extraction
        self.msst_blocks = nn.ModuleList([
            self.create_cnn_block(num_channels, 32, kernel_size=k)
            for k in [3, 5, 7]
        ])
        
        # BLOCK 2: Transformer Encoder
        self.transformer = nn.TransformerEncoder(
            nn.TransformerEncoderLayer(
                d_model=96,
                nhead=4,
                dim_feedforward=256,
                dropout=0.5
            ),
            num_layers=2
        )
        
        # BLOCK 3: Classification Head
        self.classifier = nn.Sequential(
            nn.Linear(96, 32),
            nn.ReLU(),
            nn.Dropout(0.5),
            nn.Linear(32, num_classes)
        )
\end{lstlisting}

\subsection{Step 5: Training Process}

\subsubsection{Non-Technical Explanation}

Training is like \textbf{teaching a child}:

\begin{tcolorbox}[colback=yellow!10,colframe=orange!50!black,title=Training Analogy]
\begin{enumerate}
    \item Show 50 examples of LEFT HAND brain signals
    \item Model guesses: LEFT HAND or RIGHT HAND
    \item If wrong: No, that was LEFT HAND! Learn from this!
    \item Adjust model's brain (weights) slightly
    \item Repeat thousands of times
    \item Eventually: Model becomes expert
\end{enumerate}
\end{tcolorbox}

\subsubsection{Technical Explanation}

\begin{lstlisting}[caption={Training Loop with Quality Checks}]
import torch.optim as optim

optimizer = optim.Adam(model.parameters(), lr=0.001)
loss_fn = nn.CrossEntropyLoss()

for epoch in range(150):
    for batch_data, batch_labels in train_loader:
        
        # YOUR QUALITY CHECK
        if torch.isnan(batch_data).any():
            print(f"ERROR: NaN at epoch {epoch}")
            continue
        
        predictions = model(batch_data)
        loss = loss_fn(predictions, batch_labels)
        
        optimizer.zero_grad()
        loss.backward()
        optimizer.step()
    
    print(f"Epoch {epoch}: Loss={loss:.4f}")
\end{lstlisting}

\subsection{Step 6: Testing and Evaluation}

\subsubsection{Non-Technical Explanation}

After training, we test on \textbf{NEW data the model has never seen}:

\begin{tcolorbox}[colback=green!10,colframe=green!50!black,title=Testing Process]
\begin{enumerate}
    \item Show model 30\% of signals it has never seen
    \item Model makes predictions
    \item Compare predictions with actual labels
    \item Calculate accuracy
    \item Expected: 90-95\% accuracy
\end{enumerate}
\end{tcolorbox}

\subsubsection{Technical Explanation}

\begin{lstlisting}[caption={Evaluation with Quality Metrics}]
import numpy as np
from sklearn.metrics import accuracy_score, confusion_matrix

model.eval()
all_predictions = []
all_labels = []

with torch.no_grad():
    for test_data, test_labels in test_loader:
        
        if test_data.shape[0] == 0:
            print("ERROR: Empty batch!")
            continue
        
        predictions = model(test_data)
        preds = torch.argmax(predictions, dim=1).cpu().numpy()
        labels = test_labels.cpu().numpy()
        
        all_predictions.extend(preds)
        all_labels.extend(labels)

all_predictions = np.array(all_predictions)
all_labels = np.array(all_labels)

print(f"Predictions with NULL: {np.isnan(all_predictions).sum()}")
print(f"Labels with NULL: {np.isnan(all_labels).sum()}")

accuracy = accuracy_score(all_labels, all_predictions)
print(f"Accuracy: {accuracy:.4f}")
\end{lstlisting}

% ============================================================================
\section{Your Data Quality Validation Checklist}
% ============================================================================

\subsection{Complete Data Validation Script}

\begin{lstlisting}[caption={Complete Validation (Your Main Task)}]
import numpy as np
import pandas as pd
import glob
import os

def validate_data_quality(data_dir):
    print("="*60)
    print("DATA QUALITY VALIDATION REPORT")
    print("="*60)
    
    files = glob.glob(os.path.join(data_dir, "*.npz"))
    total_issues = 0
    
    for file in files:
        print(f"\nValidating: {os.path.basename(file)}")
        print("-" * 40)
        
        try:
            data = np.load(file)
            signals = data['data']
            labels = data['labels']
            
            # CHECK 1: NULL Values
            print("CHECK 1: NULL Values")
            nan_signals = np.isnan(signals).sum()
            nan_labels = np.isnan(labels).sum()
            inf_signals = np.isinf(signals).sum()
            
            if nan_signals > 0:
                print(f"  WARNING: {nan_signals} NaN values in signals!")
                total_issues += 1
            else:
                print(f"  OK: No NaN in signals")
            
            if nan_labels > 0:
                print(f"  WARNING: {nan_labels} NaN in labels!")
                total_issues += 1
            else:
                print(f"  OK: No NaN in labels")
            
            # CHECK 2: Shape Consistency
            print("\nCHECK 2: Shape Consistency")
            if signals.shape != (signals.shape[0], 22, 875):
                print(f"  WARNING: Unexpected shape {signals.shape}")
                total_issues += 1
            else:
                print(f"  OK: Shape is {signals.shape}")
            
            # CHECK 3: Labels Validity
            print("\nCHECK 3: Labels Validity")
            unique_labels = np.unique(labels)
            print(f"  Unique labels: {unique_labels}")
            
            if not set(unique_labels).issubset({0, 1, 2, 3}):
                print(f"  WARNING: Invalid label values!")
                total_issues += 1
            else:
                print(f"  OK: All labels valid")
            
            # CHECK 4: Class Balance
            print("\nCHECK 4: Class Balance")
            unique, counts = np.unique(labels, return_counts=True)
            for label, count in zip(unique, counts):
                print(f"  Class {int(label)}: {count} samples")
            
            # CHECK 5: Signal Statistics
            print("\nCHECK 5: Signal Statistics")
            print(f"  Mean: {signals.mean():.4f} (should be 0)")
            print(f"  Std:  {signals.std():.4f} (should be 1)")
            print(f"  Min:  {signals.min():.4f}")
            print(f"  Max:  {signals.max():.4f}")
            
        except Exception as e:
            print(f"  ERROR loading file: {e}")
            total_issues += 1
    
    print(f"\n{'='*60}")
    print(f"TOTAL ISSUES FOUND: {total_issues}")
    status = "PASS" if total_issues == 0 else "REVIEW NEEDED"
    print(f"Status: {status}")
    print(f"{'='*60}\n")
    
    return total_issues == 0

if __name__ == "__main__":
    data_dir = "path/to/preprocessed/data"
    is_valid = validate_data_quality(data_dir)
    
    if is_valid:
        print("Data is ready for training!")
    else:
        print("Data needs cleaning before training!")
\end{lstlisting}

% ============================================================================
\section{Interview Questions and Answers}
% ============================================================================

\subsection{Question 1: What is the main purpose of this project?}

\textbf{Answer:}

The MSVTNet project is designed to \textbf{recognize motor imagery} (brain signals representing thought about movement) and classify them into 4 categories: Left Hand, Right Hand, Feet, or Tongue. 

\textbf{Application:} Brain-Computer Interface (BCI) for helping paralyzed patients control prosthetic limbs or wheelchairs using only their thoughts.

\textbf{Key Achievement:} 94 percent plus accuracy on real EEG data.

---

\subsection{Question 2: What is your role in this project?}

\textbf{Answer:}

My role is \textbf{Data Quality and Validation Specialist}. My responsibilities include:

\begin{enumerate}
    \item Checking for NULL/NaN values in brain signal data using NumPy and Pandas
    \item Validating data integrity across all datasets
    \item Identifying anomalies (outliers, incorrect values)
    \item Ensuring consistency before model training begins
    \item Generating quality reports for stakeholders
\end{enumerate}

I use Python libraries like \textbf{NumPy} for array operations and \textbf{Pandas} for data frame analysis.

---

\subsection{Question 3: How do you check for NULL values?}

\textbf{Answer:}

\begin{lstlisting}
import numpy as np
import pandas as pd

# Load data
data = np.load('preprocessed_data.npz')
signals = data['data']
labels = data['labels']

# Method 1: NumPy (for arrays)
nan_count = np.isnan(signals).sum()
print(f"Total NaN values: {nan_count}")

# Method 2: Pandas
df = pd.DataFrame(signals.reshape(signals.shape[0], -1))
print(f"Null per column:\n{df.isnull().sum()}")

# Method 3: Check for infinite values
inf_count = np.isinf(signals).sum()
print(f"Infinite values: {inf_count}")

# Complete validation
if nan_count == 0 and inf_count == 0:
    print("Data is clean!")
else:
    print("Data contains anomalies")
\end{lstlisting}

---

\subsection{Question 4: What are the three main blocks of MSVTNet?}

\textbf{Answer:}

\begin{enumerate}
    \item \textbf{MSST Block (Multi-Scale Spatio-Temporal)}
    \begin{itemize}
        \item Extracts patterns at 3 different time scales
        \item Uses CNN with kernel sizes: 3, 5, 7
        \item Captures both fast and slow patterns in brain signals
    \end{itemize}
    
    \item \textbf{CSGT Block (Cross-Scale Global Temporal Encoder)}
    \begin{itemize}
        \item Uses Transformer with multi-head attention
        \item Learns relationships between patterns from different scales
        \item Combines 3 feature maps into 1 unified representation
    \end{itemize}
    
    \item \textbf{Classification Head}
    \begin{itemize}
        \item Fully connected layers plus ReLU activation
        \item Makes final prediction
        \item Output: 4 class probabilities
    \end{itemize}
\end{enumerate}

---

\subsection{Question 5: What preprocessing steps are applied to raw data?}

\textbf{Answer:}

\begin{enumerate}
    \item \textbf{Bandpass Filter (4-40 Hz)}
    \begin{itemize}
        \item Removes electrical noise outside motor cortex frequency range
        \item Keeps only relevant brain signals
    \end{itemize}
    
    \item \textbf{Common Average Reference (CAR)}
    \begin{itemize}
        \item Subtracts mean of all channels from each channel
        \item Removes common environmental noise
    \end{itemize}
    
    \item \textbf{Standardization (Z-score Normalization)}
    \begin{itemize}
        \item Converts to mean = 0, standard deviation = 1
        \item Makes signals comparable across subjects
    \end{itemize}
    
    \item \textbf{Epoch Extraction}
    \begin{itemize}
        \item Extracts 3.5-second windows (875 samples at 250 Hz)
        \item Labels each window with movement class
    \end{itemize}
\end{enumerate}

---

\subsection{Question 6: How do you validate training data before feeding to model?}

\textbf{Answer:}

\begin{lstlisting}
def validate_before_training(train_loader):
    for batch_idx, (data, labels) in enumerate(train_loader):
        
        # Check 1: No NULL values
        assert not torch.isnan(data).any(), \
            f"NaN found in batch {batch_idx}"
        
        # Check 2: Correct shape
        assert data.shape[1:] == (22, 875), \
            f"Unexpected shape: {data.shape}"
        
        # Check 3: Labels in valid range
        assert labels.min() >= 0 and labels.max() < 4, \
            f"Invalid labels: {labels.unique()}"
        
        # Check 4: No empty batches
        assert data.shape[0] > 0, "Empty batch!"
        
        if batch_idx == 0:
            print("First batch validation passed!")
            print(f"  Data: {data.shape}, Labels: {labels.shape}")
\end{lstlisting}

---

\subsection{Question 7: What is Cross-Entropy Loss?}

\textbf{Answer:}

\textbf{Simple Explanation:} A mathematical formula that measures how wrong the model's prediction is.

\begin{itemize}
    \item If model predicts LEFT HAND correctly: Loss is approximately 0 (very good)
    \item If model predicts TONGUE but actual is LEFT HAND: Loss is approximately 2.5 (very bad)
\end{itemize}

\textbf{Mathematical Formula:}
\begin{equation}
L = -\sum_{i=1}^{C} y_i \log(\hat{y}_i)
\end{equation}

Where:
\begin{itemize}
    \item $y_i$ = actual label (1 if correct class, 0 otherwise)
    \item $\hat{y}_i$ = predicted probability for class $i$
    \item $C$ = number of classes (4 in our case)
\end{itemize}

\textbf{Goal during training:} Minimize this loss value.

---

\subsection{Question 8: What is Backpropagation?}

\textbf{Answer:}

\textbf{Simple Analogy:} Like throwing darts at a target:

\begin{enumerate}
    \item You throw dart --- It misses
    \item You measure: Missed 10cm to the right
    \item Next throw: Adjust aim 10cm to the left
    \item Repeat: Dart gets closer each time
\end{enumerate}

\textbf{In AI:}

\begin{enumerate}
    \item Model makes prediction --- Calculates loss
    \item Works backwards through all layers
    \item Calculates: Which weights caused this error?
    \item Adjusts all weights slightly in the right direction
    \item Repeat: Model gets better each epoch
\end{enumerate}

---

\subsection{Question 9: What does Dropout do?}

\textbf{Answer:}

\textbf{Purpose:} Prevent overfitting (memorization instead of learning).

\textbf{How it works:}

During training:
\begin{itemize}
    \item Randomly turn off 50 percent of neurons
    \item Network must learn without relying on any single neuron
    \item Makes model more robust
\end{itemize}

During testing:
\begin{itemize}
    \item All neurons are ON
    \item Model uses everything it learned
\end{itemize}

---

\subsection{Question 10: What is Early Stopping?}

\textbf{Answer:}

\textbf{Purpose:} Stop training when model stops improving to prevent overfitting.

\textbf{How it works:}

\begin{enumerate}
    \item Monitor validation accuracy after each epoch
    \item If accuracy improves: Continue training
    \item If accuracy gets worse for N consecutive epochs: STOP
    \item Use the model from best epoch
\end{enumerate}

% ============================================================================
\section{Key Algorithms Used}
% ============================================================================

\subsection{1. Convolutional Neural Network (CNN)}
Purpose: Extract patterns from sequential data (brain signals)

\subsection{2. Transformer (Self-Attention)}
Purpose: Learn relationships between time steps and features

\subsection{3. Multi-Head Attention}
Purpose: Focus on different types of patterns simultaneously

\subsection{4. Backpropagation}
Purpose: Learn by adjusting weights based on prediction errors

\subsection{5. ADAM Optimizer}
Purpose: Update weights with adaptive learning rates

\subsection{6. Dropout}
Purpose: Prevent overfitting by random neuron deactivation

\subsection{7. Batch Normalization}
Purpose: Stabilize training and normalize data within batches

\subsection{8. Cross-Entropy Loss}
Purpose: Measure how wrong model predictions are

% ============================================================================
\section{Summary and Key Takeaways}
% ============================================================================

\subsection{Project Overview}
\begin{itemize}
    \item \textbf{Goal:} Classify EEG brain signals into 4 movement classes
    \item \textbf{Accuracy:} 94 percent plus on real data
    \item \textbf{Application:} Brain-Computer Interface for rehabilitation
    \item \textbf{Technology:} PyTorch, Transformers, CNNs
\end{itemize}

\subsection{Your Role}
\begin{itemize}
    \item \textbf{Data Quality Specialist}
    \item Check for NULL/NaN values using NumPy and Pandas
    \item Validate data integrity before training
    \item Generate quality reports
    \item Ensure consistent data shapes and types
\end{itemize}

\subsection{Training Pipeline}
\begin{enumerate}
    \item \textbf{Data Collection:} Record EEG signals from subjects
    \item \textbf{Preprocessing:} Filter, normalize, and clean data
    \item \textbf{Data Validation:} Your role - check quality
    \item \textbf{Model Training:} Feed data, calculate loss, backpropagate
    \item \textbf{Testing:} Evaluate on unseen data
    \item \textbf{Deployment:} Use model for BCI applications
\end{enumerate}

\subsection{Key Metrics to Remember}
\begin{table}[h]
\centering
\begin{tabular}{|l|l|}
\hline
\textbf{Metric} & \textbf{Expected Value} \\
\hline
Accuracy & 90-95 percent \\
Loss & less than 0.5 \\
NaN count & 0 \\
Signal mean & approximately 0 \\
Signal std & approximately 1 \\
\hline
\end{tabular}
\end{table}

\section{Quick Reference: Daily Tasks}

\begin{lstlisting}[caption={Your Daily Task - Check NULL Values}]
import numpy as np
import pandas as pd

# Load data
data = np.load('S01_session_1_preprocessed.npz')
signals = data['data']
labels = data['labels']

# Quick checks
print(f"NaN in signals: {np.isnan(signals).sum()}")
print(f"NaN in labels: {np.isnan(labels).sum()}")
print(f"Data shape: {signals.shape}")
print(f"Label distribution: {np.bincount(labels)}")

# If using Pandas
df = pd.DataFrame({
    'signal_mean': signals.mean(axis=(1,2)),
    'signal_std': signals.std(axis=(1,2)),
    'label': labels
})
print(df.isnull().sum())
\end{lstlisting}

% ============================================================================
\end{document}
